% 06-examples-simulations
\chapter{Examples and simulations}\label{ch:examples}

\Cref{ch:turing} turned the question ``can diffusion destabilize?'' into conditions
on $2\times2$ matrices and a band of unstable modes. In this chapter we make these
results concrete. First we draw, for our running examples, the regions in parameter
space where diffusion-driven instability (DDI) is possible. Then we set up a simple
numerical method and use it to solve the full nonlinear system. The simulations
confirm the linear predictions of \cref{ch:turing} in the early phase, with
surprising accuracy, and they show what linear theory cannot tell us: how large the
pattern becomes, and which of several competing modes wins. All numbers quoted below
were produced by short Python scripts that accompany these notes.

\section{Turing spaces}\label{sec:turing-spaces}

Kinetics usually depend on parameters, such as $a,b$ in our running examples. The
Turing conditions (T1)--(T4) of \eqref{eq:turing-conditions} involve the diffusion
coefficients only through the ratio $\delta=d_v/d_u$. For a fixed value of $\delta$
we call the set of kinetic parameters for which (T1)--(T4) hold the \emph{Turing
space}. It does not depend on the interval: on $(0,L)$ an instability occurs only if,
in addition, some $\lambda_k$ lies in the band \eqref{eq:band}, which is automatic
for large $L$ by \cref{cor:domain-size}.

\begin{example}[Schnakenberg]\label{ex:examples-schnakenberg-space}
By \cref{ex:schnakenberg}, $f_u=(b-a)/(a+b)$, $f_v=(a+b)^2$, $g_u=-2b/(a+b)$,
$g_v=-(a+b)^2$ and $\det A=(a+b)^2$. So (T2) always holds, and (T1) reads
$(b-a)/(a+b)<(a+b)^2$. Since $g_v<0$, \cref{cor:turing-consequences}(a) shows that
DDI requires $f_u>0$, i.e.\ $a<b$. Then \cref{thm:critical-ratio} applies with
$-f_vg_u=2b(a+b)$, and the Turing space for $\delta$ is the set of $a,b>0$ with
\begin{equation}\label{eq:examples-schnakenberg-space}
  a<b,\qquad b-a<(a+b)^3,\qquad
  \delta>\delta_c(a,b)=\frac{(a+b)^2\big((a+b)+\sqrt{2b(a+b)}\,\big)^2}{(b-a)^2}.
\end{equation}
For $(a,b)=(0.1,0.9)$ we recover $\delta_c=(1+\sqrt{1.8})^2/0.64\approx8.57$
(\cref{ex:schnakenberg-turing}).
\end{example}

\begin{figure}[t]
  \centering
  \includegraphics[width=\textwidth]{ch06-turing-space}
  \caption{Turing spaces. (a) Schnakenberg kinetics in the $(a,b)$ plane: shaded,
  the parameters with DDI for $\delta=20$, bounded by the curve $\delta_c(a,b)=20$
  and by the region where the kinetics are unstable ((T1) fails, grey); dashed, the
  curve $\delta_c(a,b)=10$. The dot is the standard point $(0.1,0.9)$. (b)
  Gierer--Meinhardt kinetics with $a=0$ in the $(b,\delta)$ plane: DDI occurs above
  the curve $\delta=(1+\sqrt2)^2/b$ and left of $b=1$. Filled dot: the standard
  parameters $b=0.5$, $\delta=30$; open dot: $\delta_c\approx11.66$ for $b=0.5$.}
  \label{fig:ch06-turing-space}
\end{figure}

\Cref{fig:ch06-turing-space}(a) shows this region. For $a=0.1$, condition (T1)
requires $b>0.779$, and the Turing space consists of the values $0.779<b<1.002$ for
$\delta=10$ and $0.779<b<1.559$ for $\delta=20$. Thus the standard point is rather
close to the region where the kinetics themselves are unstable; there the reaction
oscillates in time, which was Schnakenberg's original interest \cite{Schnakenberg1979}.
The Turing space grows with $\delta$: since $\delta_c(a,b)<\infty$ whenever $a<b$,
every point with $a<b$ and (T1) belongs to it once $\delta$ is large enough. Near the
curve $\delta_c(a,b)=\delta$ the instability is weak; we will see in
\cref{sec:examples-2d} that patterns near this edge and deep inside can look quite
different.

\begin{example}[Gierer--Meinhardt, $a=0$]\label{ex:examples-gm-space}
For $a=0$, \cref{ex:gierer-meinhardt} gives $f_u=b$, $f_v=-b^2$, $g_u=2/b$, $g_v=-1$,
$\det A=b$ and $\tr A=b-1$. Hence (T2) and $f_u>0$ always hold, (T1) means $b<1$, and
$-f_vg_u=2b$. \Cref{thm:critical-ratio} gives
\begin{equation*}
  \delta_c(b)=\frac{\big(\sqrt b+\sqrt{2b}\,\big)^2}{b^2}=\frac{(1+\sqrt2)^2}{b}
  \approx\frac{5.83}{b},\qquad 0<b<1 .
\end{equation*}
The Turing space is the region above this hyperbola in \cref{fig:ch06-turing-space}(b).
The general case $a\ge0$ is \cref{exer:examples-gm-space}.
\end{example}

\section{A numerical method}\label{sec:numerics}

We discretize in space only and keep time continuous; this turns the PDE into a large
system of ODEs (the \emph{method of lines}), which we then solve by a time-stepping
scheme. Divide $(0,L)$ into $N$ cells of width $h=L/N$ (not to be confused with the
function $h(\lambda)$ of \cref{ch:turing}) with centers
\begin{equation*}
  x_j=\big(j-\tfrac12\big)h,\qquad j=1,\dots,N,
\end{equation*}
and let $u_j(t)\approx u(x_j,t)$. We replace $u_{xx}$ by the central difference
quotient and impose the Neumann condition by \emph{reflected ghost values}
$u_0:=u_1$, $u_{N+1}:=u_N$. (The ghost point $x_0=-h/2$ is the mirror image of $x_1$,
and $(u_1-u_0)/h=0$ is a central difference for $u_x(0)=0$.) This defines the
\emph{discrete Laplacian} $\Delta_h:\R^N\to\R^N$,
\begin{equation}\label{eq:discrete-laplacian}
  (\Delta_hu)_j=\frac{u_{j+1}-2u_j+u_{j-1}}{h^2},\qquad j=1,\dots,N,
\end{equation}
and the semi-discrete system $u_j'=d_u(\Delta_hu)_j+f(u_j,v_j)$,
$v_j'=d_v(\Delta_hv)_j+g(u_j,v_j)$. Remarkably, the discrete problem has exactly the
same cosine modes as the continuous one.

\begin{proposition}[Discrete cosine modes]\label{prop:examples-discrete-modes}
For $k=0,1,\dots,N-1$ the vector $\varphi^h_k=\big(\cos(k\pi x_j/L)\big)_{j=1}^N$
satisfies
\begin{equation*}
  \Delta_h\varphi^h_k=-\lambda^h_k\,\varphi^h_k,\qquad
  \lambda^h_k=\frac4{h^2}\sin^2\Big(\frac{k\pi h}{2L}\Big).
\end{equation*}
These $N$ vectors form an orthogonal basis of $\R^N$, and
$\lambda^h_k\to\lambda_k=(k\pi/L)^2$ as $h\to0$ for each fixed $k$.
\end{proposition}

\begin{proof}
Put $\theta=k\pi h/L$, so that $\varphi^h_{k,j}=\cos\big((j-\tfrac12)\theta\big)$.
From $\cos(\alpha+\theta)+\cos(\alpha-\theta)=2\cos\alpha\cos\theta$ with
$\alpha=(j-\frac12)\theta$ we get
\begin{equation*}
  \cos\big((j+\tfrac12)\theta\big)-2\cos\big((j-\tfrac12)\theta\big)
  +\cos\big((j-\tfrac32)\theta\big)
  =-2(1-\cos\theta)\cos\big((j-\tfrac12)\theta\big)
  =-4\sin^2(\theta/2)\,\cos\big((j-\tfrac12)\theta\big).
\end{equation*}
This is the claim for $2\le j\le N-1$. For $j=1$ and $j=N$ it also holds, because
the reflected ghost values agree with the cosine formula: $\cos(-\theta/2)=\cos(\theta/2)$,
and, since $N\theta=k\pi$, $\cos\big((N+\frac12)\theta\big)=(-1)^k\cos(\theta/2)
=\cos\big((N-\frac12)\theta\big)$. The matrix of $\Delta_h$ is symmetric (its rows are
$(-1,1,0,\dots)/h^2$, $(\dots,1,-2,1,\dots)/h^2$ and $(\dots,0,1,-1)/h^2$), and the
$\lambda^h_k$ are pairwise distinct because $\sin^2$ is strictly increasing on
$[0,\pi/2)$ and $k\pi h/(2L)=k\pi/(2N)\in[0,\pi/2)$. Eigenvectors of a symmetric
matrix for distinct eigenvalues are orthogonal, so we have $N$ orthogonal nonzero
vectors in $\R^N$. Finally, with $s=k\pi h/(2L)$ we have
$\lambda^h_k=\lambda_k(\sin s/s)^2\to\lambda_k$ as $h\to0$.
\end{proof}

So everything we did in \cref{ch:turing} has an exact discrete counterpart: the
semi-discrete linearized system decouples into $2\times2$ systems with the mode
matrices $A-\lambda^h_kD$, $k=0,\dots,N-1$. For the modes that matter, $\lambda^h_k$
is very close to $\lambda_k$: for $L=40$, $h=0.1$ and $k=6$ they differ by $0.02\%$.

\paragraph{Explicit Euler and the step restriction.}
The simplest time-stepping scheme is the explicit Euler method with step $\Delta t$:
$u^{n+1}_j=u^n_j+\Delta t\,\big(d_u(\Delta_hu^n)_j+f(u^n_j,v^n_j)\big)$, and similarly
for $v$. The following listing is a complete implementation for the Schnakenberg
example (Python with \texttt{numpy}); the arrays \texttt{u}, \texttt{v} hold
$(u_j)$, $(v_j)$, and \texttt{lap} is \eqref{eq:discrete-laplacian}.
{\small
\begin{verbatim}
import numpy as np

a, b, du, dv = 0.1, 0.9, 1.0, 20.0      # Schnakenberg, standard values
L, N, T = 40.0, 200, 200.0
h = L / N
x = (np.arange(1, N + 1) - 0.5) * h      # cell centres
dt = 0.9 * h**2 / (2 * max(du, dv))      # step restriction, with a margin

def lap(w):                              # discrete Laplacian, Neumann by reflection
    wg = np.concatenate(([w[0]], w, [w[-1]]))
    return (wg[2:] - 2 * wg[1:-1] + wg[:-2]) / h**2

rng = np.random.default_rng(0)
u = a + b + 1e-3 * rng.uniform(-1, 1, N)          # constant state plus noise
v = b / (a + b)**2 + 1e-3 * rng.uniform(-1, 1, N)
for n in range(int(T / dt)):
    f = a - u + u**2 * v
    g = b - u**2 * v
    u, v = u + dt * (du * lap(u) + f), v + dt * (dv * lap(v) + g)
\end{verbatim}
}
\noindent It runs in a few seconds and ends with three peaks of $u$, at
$x\approx6.7$, $19.9$, $33.3$. The choice of \texttt{dt} is not arbitrary.

\begin{proposition}[Step restriction]\label{prop:examples-step-restriction}
Let $d>0$ and consider explicit Euler for the semi-discrete heat equation
$w'=d\,\Delta_hw$, i.e.\ $w^{n+1}=(I+d\,\Delta t\,\Delta_h)w^n$. If
$\Delta t\le h^2/(2d)$, then $\abs{w^n}\le\abs{w^0}$ for all $n$ (Euclidean norm in
$\R^N$). If $\Delta t>h^2/\big(2d\cos^2(\pi/(2N))\big)$, there are $w^0$ with
$\abs{w^n}\to\infty$.
\end{proposition}

\begin{proof}
Expand $w^0=\sum_kc_k\varphi^h_k$. By \cref{prop:examples-discrete-modes},
$w^n=\sum_k m_k^n c_k\varphi^h_k$ with the multipliers
$m_k=1-\frac{4d\Delta t}{h^2}\sin^2\big(\frac{k\pi}{2N}\big)$. If
$\Delta t\le h^2/(2d)$, then $0\le\frac{4d\Delta t}{h^2}\sin^2(\cdot)\le2$, so
$m_k\in[-1,1]$, and orthogonality gives
$\abs{w^n}^2=\sum_k m_k^{2n}c_k^2\abs{\varphi^h_k}^2\le\abs{w^0}^2$. For $k=N-1$ we have
$\sin^2\big(\frac{(N-1)\pi}{2N}\big)=\cos^2\big(\frac{\pi}{2N}\big)$; under the second
assumption $m_{N-1}<-1$, and $w^0=\varphi^h_{N-1}$ gives $\abs{w^n}=\abs{m_{N-1}}^n\abs{w^0}\to\infty$.
\end{proof}

For the system we therefore need $\Delta t\le h^2/(2\max\{d_u,d_v\})$; the reactions
add only the mild requirement that $\Delta t$ be small compared with the time scale
of the kinetics. If the restriction is violated, the mode $k=N-1$, which alternates
in sign from cell to cell, is amplified at every step and quickly destroys the
solution (\cref{exer:examples-step-restriction}).

\begin{remark}[Numerical versus Turing instability]\label{rem:examples-numerical-instability}
Both instabilities make a constant state blow up into a spatial pattern, but they are
very different. The numerical instability lives at the grid scale, depends on $h$
and $\Delta t$, and disappears when $\Delta t$ is reduced. The Turing instability
lives in a fixed band of wavelengths, independent of $h$ and $\Delta t$; refining the
discretization only makes it more accurate.
\end{remark}

\paragraph{A semi-implicit scheme.}
The restriction $\Delta t\le h^2/(2d_v)$ is severe: for $h=0.1$ and $d_v=20$ it gives
$\Delta t\le2.5\cdot10^{-4}$, and a run up to $t=600$ needs more than two million
steps. The remedy is to treat diffusion implicitly and the reactions explicitly
(an \emph{IMEX} scheme):
\begin{equation*}
  (I-\Delta t\,d_u\Delta_h)\,u^{n+1}=u^n+\Delta t\,f(u^n,v^n),\qquad
  (I-\Delta t\,d_v\Delta_h)\,v^{n+1}=v^n+\Delta t\,g(u^n,v^n).
\end{equation*}
Now diffusion multiplies mode $k$ by $1/(1+d\,\Delta t\,\lambda^h_k)\in(0,1]$, for
every $\Delta t>0$. The linear systems are cheap to solve, since they are diagonal in
the basis $\varphi^h_k$ (the transformation to this basis is the \emph{discrete cosine
transform}, available as \texttt{scipy.fft.dct}). A stationary solution of the scheme,
$u^{n+1}=u^n$, $v^{n+1}=v^n$, solves $d_u\Delta_hu+f=0$, $d_v\Delta_hv+g=0$ exactly, for
any $\Delta t$. All simulations below use this IMEX scheme with $h=0.1$ and
$\Delta t=0.01$ in one dimension.

\section{One dimension: watching the instability}\label{sec:examples-1d}

We simulate the standard example of \cref{ex:schnakenberg-turing}: Schnakenberg
kinetics with $a=0.1$, $b=0.9$, $d_u=1$, $d_v=20$ on $(0,40)$, where exactly the
modes $k=4,\dots,10$ are unstable and $k=6$ grows fastest, with $\mu\approx0.262$. The
initial data are $(\ub,\vb)=(1,0.9)$ plus independent random numbers, uniform in
$[-10^{-3},10^{-3}]$, at every grid point (with a fixed random seed). The result is
shown in \cref{fig:ch06-1d-evolution}.

\begin{figure}[t]
  \centering
  \includegraphics[width=\textwidth]{ch06-1d-evolution}
  \caption{Schnakenberg kinetics, $a=0.1$, $b=0.9$, $d_u=1$, $d_v=20$, $L=40$,
  initial data $(\ub,\vb)$ plus noise of size $10^{-3}$ (IMEX scheme, $h=0.1$,
  $\Delta t=0.01$). (a) $u(x,t)$ for $0\le t\le150$: nothing visible happens until
  $t\approx30$, then three peaks appear and slowly move to their final positions.
  (b) Cosine coefficients $\abs{\hat u_k(t)}$ (log scale) for four unstable modes;
  black dashed lines have the slopes $\mu_k$ predicted by linear theory. (c) The
  profile at $t=25$, scaled to maximum $1$, the solution of the linearized problem
  with the same initial data, and the fastest mode. (d) The stationary pattern: $u$
  and $v$ are out of phase; dotted lines show $\ub=1$ and $\vb=0.9$.}
  \label{fig:ch06-1d-evolution}
\end{figure}

\paragraph{The linear phase.}
By \cref{thm:modal-decomposition}, as long as the perturbation is small each cosine
coefficient should evolve by its own $2\times2$ system and, after a short transient,
grow or decay like $e^{\mu_kt}$. We compute from the simulation the discrete cosine
coefficients $\hat u_k(t)=\frac2N\sum_j(u_j(t)-\ub)\cos(k\pi x_j/L)$ (the midpoint rule
for \eqref{eq:cosine-coefficients}). In \cref{fig:ch06-1d-evolution}(b) they are
straight lines on a logarithmic scale up to $t\approx30$. A least-squares fit of
$\log\abs{\hat u_k(t)}$ on $10\le t\le20$ gives the rates in \cref{tab:examples-rates}.

\begin{table}[ht]
  \centering\small
  \caption{Growth rates of the unstable modes on $(0,40)$: eigenvalue $\mu_k$ of $A_k$
  and rate fitted from the simulation of \cref{fig:ch06-1d-evolution} on $10\le t\le20$.}
  \label{tab:examples-rates}
  \begin{tabular}{lccccccc}
    \toprule
    $k$ & 4 & 5 & 6 & 7 & 8 & 9 & 10\\
    \midrule
    linear theory $\mu_k$ & 0.1194 & 0.2284 & 0.2623 & 0.2510 & 0.2075 & 0.1386 & 0.0487\\
    simulation & 0.1186 & 0.2285 & 0.2621 & 0.2505 & 0.2070 & 0.1383 & 0.0429\\
    \bottomrule
  \end{tabular}
\end{table}

The agreement is excellent, except for the barely unstable mode $k=10$. The solution
of the linearized problem gives $0.0489$ for this mode, so the discrepancy is a
nonlinear effect: the quadratic terms of $f,g$ multiply growing modes, and products
such as $\cos(4\pi x/L)\cos(6\pi x/L)=\frac12\big(\cos(10\pi x/L)+\cos(2\pi x/L)\big)$
feed other modes, among them stable ones. At $t=25$ the perturbation is still small
($\max\abs{u-\ub}=0.041$) and differs from the linearized solution by $1.4\%$ of its
size (\cref{fig:ch06-1d-evolution}(c)). Note, however, that it is \emph{not} yet the
pure mode $\cos(6\pi x/L)$: mode 6 carries only $49\%$ of $\norm{u-\ub}^2$, because its
neighbors grow almost as fast.

\paragraph{Saturation and selection.}
At $t\approx40$ the perturbation reaches size one, the growth stops, and the peaks
then drift slowly to $x\approx6.65$, $19.95$, $33.35$, the maxima of $-\cos(6\pi x/40)$
(up to the grid). From $t\approx200$ on the solution is stationary up to $10^{-2}$,
and at $t=600$ it changes by less than $10^{-7}$ in ten time units. This run thus
confirms the linear prediction: the final pattern has the shape of the fastest
mode $k=6$. But this is not the whole story. We repeated the run with $50$ different
random seeds: the final pattern was of mode $6$ in $31$ runs, of mode $7$ (peaks at
$x\approx5.75$, $17.15$, $28.55$, $39.95$) in $18$ runs, and of mode $5$ once. Both
mode-$6$ and mode-$7$ patterns are stable stationary solutions. The explanation is
that $\mu_6-\mu_7=0.011$: during the roughly $35$ time units of linear growth, mode $6$
gains only a factor $e^{0.011\cdot35}\approx1.5$ over mode $7$, which the random initial
amplitudes can easily outweigh.

\begin{keyidea}
Linear theory correctly predicts the early phase: which modes grow, how fast, and
hence the length scale of the pattern. It does not predict the final amplitude, and
when several modes grow at nearly the same rate it does not decide which one wins.
These are nonlinear questions, and the answer may depend on the initial data.
\end{keyidea}

\paragraph{In phase and out of phase.}
\Cref{rem:turing-phase} predicts that the unstable modes of substrate-depletion
kinetics ($g_u<0$) have $u$ and $v$ out of phase, and those of activator--inhibitor
kinetics ($g_u>0$) in phase. \Cref{fig:ch06-phase-relation} shows that this sign
relation persists in the final nonlinear patterns. For the Gierer--Meinhardt example
we take $a=0$, $b=0.5$, $d_u=1$, $d_v=30$ on $(0,20)$, where the modes $k=2,3,4$ are
unstable and $k=2$ is the fastest. The quotient $u^2/v$ requires $v>0$; the exact
solution keeps $v$ positive, since $v_t\ge d_vv_{xx}-v$, and so does the IMEX scheme
(\cref{exer:examples-imex}). In the run shown, $v\ge2.75$ throughout.

\begin{figure}[t]
  \centering
  \includegraphics[width=\textwidth]{ch06-phase-relation}
  \caption{Stationary patterns, plotted as $(u-\ub)/\max\abs{u-\ub}$ and
  $(v-\vb)/\max\abs{v-\vb}$. (a) Schnakenberg ($a=0.1$, $b=0.9$, $d_v=20$, $L=40$),
  substrate depletion: the peaks of $u$ sit at the troughs of $v$. (b)
  Gierer--Meinhardt ($a=0$, $b=0.5$, $d_v=30$, $L=20$), activator--inhibitor: the
  peaks coincide. Both with $d_u=1$, random initial data, IMEX scheme up to $t=600$.}
  \label{fig:ch06-phase-relation}
\end{figure}

\paragraph{The size of the domain.}
\Cref{cor:domain-size} says which modes are unstable on $(0,L)$: the integers $k$
with $L\sqrt{\lambda_-}/\pi<k<L\sqrt{\lambda_+}/\pi$. In the $(L,k)$ plane these are the
lattice points in a wedge (\cref{fig:ch06-domain-size}(a)). For $L<\pi/\sqrt{\lambda_+}
\approx3.82$ there are none, and the simulation with $L=3$ indeed returns to the
constant state. For $L=10$ the modes $k=1,2$ are unstable and the final pattern is of
mode $2$, a single peak in the middle. For $L=80$ the modes $k=7,\dots,20$ are unstable,
$k=12$ and $k=13$ grow almost equally fast ($\mu\approx0.262$ and $0.261$), and our run
ends with mode $12$. In all cases with a pattern, the peaks have about the same
height and spacing, close to the wavelength $2\pi/\sqrt{\lambda^*}\approx12.9$ of the
fastest-growing wavenumber of the dispersion relation ($\lambda^*\approx0.237$). The
length scale is intrinsic, as predicted in \cref{sec:turing-criterion}; the domain
only decides how many peaks fit.

\begin{figure}[t]
  \centering
  \includegraphics[width=\textwidth]{ch06-domain-size}
  \caption{Role of the domain size, Schnakenberg with $a=0.1$, $b=0.9$, $d_u=1$,
  $d_v=20$. (a) Mode $k$ is unstable on $(0,L)$ if and only if $(L,k)$ lies in the
  shaded wedge between the lines $k=L\sqrt{\lambda_\pm}/\pi$ (\cref{cor:domain-size});
  the dots are the unstable modes for $L=3,10,40,80$ (none for $L=3$). The wedge
  begins at $L_{\min}=\pi/\sqrt{\lambda_+}\approx3.82$; the dotted line marks the
  fastest-growing wavenumber. (b) Final $u$ for these four lengths, on a common
  $x$-axis (IMEX scheme up to $t=600$, for $L=80$ up to $t=2000$).}
  \label{fig:ch06-domain-size}
\end{figure}

\section{Two dimensions}\label{sec:examples-2d}

On the square $(0,L)^2$ with Neumann conditions the same method works: the
five-point Laplacian with reflected ghost values has the eigenvectors
$\cos(m\pi x_i/L)\cos(n\pi y_j/L)$, by \cref{prop:examples-discrete-modes} applied in
each direction. The continuous Neumann eigenfunctions are
$\cos(m\pi x/L)\cos(n\pi y/L)$ with eigenvalues $(\pi/L)^2(m^2+n^2)$
(\cref{ex:rectangle-eigen}), and the linear theory of \cref{ch:turing} carries over
(\cref{rem:general-domain}): mode $(m,n)$ grows if and only if its eigenvalue lies in
the band. For the Schnakenberg standard example on $(0,80)^2$ this holds for $320$
pairs $(m,n)$, the lattice points in a quarter annulus. Many of them grow at nearly
the same rate, and some at exactly the same rate: $(m,n)$ and $(n,m)$, or
$(1,7)$ and $(5,5)$, since $1^2+7^2=5^2+5^2$. Linear theory fixes the wavelength, but
not the geometry: a single mode $\cos(m\pi x/L)$ is a pattern of stripes, while a sum
of modes with different directions produces spots.

\begin{figure}[t]
  \centering
  \includegraphics[width=\textwidth]{ch06-2d-patterns}
  \caption{Schnakenberg kinetics on the square $(0,80)^2$ with Neumann conditions,
  $a=0.1$, $d_u=1$, $d_v=20$, random initial data near $(\ub,\vb)$, state at
  $t=2000$ ($128\times128$ cells, IMEX scheme with $\Delta t=0.05$). (a), (b) $b=0.9$:
  spots of $u$, located at the minima of $v$. (c) $b=1.5$, close to the edge of the
  Turing space ($\delta_c\approx18.8$): stripes.}
  \label{fig:ch06-2d-patterns}
\end{figure}

\Cref{fig:ch06-2d-patterns} shows two simulations, with $128\times128$ cells and the
IMEX scheme with $\Delta t=0.05$. For $b=0.9$ we obtain about fifty spots of $u$, at
the minima of $v$ (out of phase, as in one dimension). For $b=1.5$ we obtain stripes.
We found this value by a small search: one run each for $a=0.1$, $\delta=20$ and
$b=0.8$, $0.9$, $1.1$, $1.3$, $1.4$, $1.5$ gave spots for $b\le1.1$, a mixture for
$b=1.3$, and stripes for $b=1.4$ and $1.5$. (Three Gierer--Meinhardt runs gave spots;
Lengyel--Epstein kinetics \eqref{eq:lengyel-epstein} with $a=10$, $b=4$ gave stripes for
$\delta=25$ and spots for $\delta=40$.) In both families the stripes appeared closer to
the threshold $\delta_c$ than the spots, but a handful of runs is not a theorem. Also,
at $t=2000$ the stripes are still slowly rearranging, while the spots have almost
settled. Which pattern is selected in two dimensions, and why, is a genuinely
nonlinear question beyond these notes; see \cref{ch:outlook} and
\cite{Murray2003}.

\section{Exercises}\label{sec:examples-exercises}

\begin{exercise}\label{exer:examples-step-restriction}
Run the listing of \cref{sec:numerics}, then replace the factor $0.9$ in \texttt{dt}
by $1.02$ and set \texttt{T = 5.0}. Plot $v$ at the end and explain what you see
using the multiplier $m_{N-1}$ from the proof of \cref{prop:examples-step-restriction}.
Why does the trouble start in $v$ and not in $u$?
\end{exercise}

\begin{exercise}\label{exer:examples-imex}
Consider the IMEX scheme of \cref{sec:numerics}.
\begin{enumerate}[label=(\alph*)]
  \item Let $c>0$ and $r\in\R^N$ with $r_j\ge0$ for all $j$. Show that the solution of
        $(I-c\Delta_h)w=r$ satisfies $w_j\ge0$ for all $j$. \emph{Hint:} look at an
        index $j$ where $w_j$ is minimal.
  \item For the Gierer--Meinhardt kinetics, show that $v^n_j>0$ for all $j$ implies
        $v^{n+1}_j>0$ for all $j$, provided $\Delta t<1$.
\end{enumerate}
\end{exercise}

\begin{exercise}\label{exer:examples-gm-space}
Consider the Gierer--Meinhardt kinetics \eqref{eq:gierer-meinhardt} with $a\ge0$.
Show that DDI is impossible for $a\ge1$, and that for $0\le a<1$ the Turing space is
given by
\begin{equation*}
  b(1-a)<1+a\qquad\text{and}\qquad
  \delta>\delta_c(a,b)=\frac{(1+a)\big(\sqrt{1+a}+\sqrt2\,\big)^2}{b\,(1-a)^2}.
\end{equation*}
Sketch the Turing space in the $(a,b)$ plane for $\delta=30$.
\end{exercise}

\begin{exercise}\label{exer:examples-l20}
Take the Schnakenberg standard example, but on $(0,20)$. Which modes are unstable, and
which grows fastest? Explain why the fastest growth rate is the same as on $(0,40)$.
Predict the final pattern, test your prediction with a few random seeds, and explain
why the outcome is more reliable than on $(0,40)$.
\end{exercise}

\begin{exercise}[$\ast$]\label{exer:examples-near-threshold}
Take the Schnakenberg kinetics with $a=0.1$, $b=0.9$, $d_u=1$ and
$d_v=\delta_c(1+\eps)$ with $\delta_c\approx8.57$. Let $\lambda_c=\delta_c^{-1/2}$
(\cref{thm:critical-ratio}) and $L=\pi/\sqrt{\lambda_c}$.
\begin{enumerate}[label=(\alph*)]
  \item Show that for small $\eps>0$ only the mode $k=1$ is unstable, and compute
        $\mu_1$ numerically for $\eps=0.01$, $0.02$, $0.05$, $0.1$. How does $\mu_1$
        depend on $\eps$?
  \item Simulate until the solution is stationary and plot the final amplitude
        $\max\abs{u-\ub}$ against $\eps$. Show that it behaves like $C\sqrt\eps$.
  \item Repeat with the Gierer--Meinhardt kinetics ($a=0$, $b=0.5$,
        $\delta_c\approx11.66$, $\lambda_c=(0.5/\delta_c)^{1/2}$). What is different?
        Relate your observations to \cref{rem:turing-bifurcation} and
        \cref{ch:outlook}.
\end{enumerate}
\end{exercise}
